%% Design Documentation: Project Smart Trigger
%% Items marked \todo{} render in red and still need team input before submission.

\documentclass[11pt, journal]{IEEEtran}
\usepackage{hyperref}
\usepackage{color}
\usepackage{float}
\usepackage[pdftex]{graphicx}
\usepackage{amsmath,amssymb,amsfonts}
\usepackage{xurl}
\usepackage[table,xcdraw]{xcolor}
\usepackage{array}
\usepackage{tikz}
\usetikzlibrary{positioning, arrows.meta, fit, backgrounds, shapes.geometric}

% Draft marker for content the team still needs to confirm or measure.
% Redefine as \newcommand{\todo}[1]{} to hide all markers for the final build.
\newcommand{\todo}[1]{\textcolor{red}{[TODO: #1]}}

\hyphenation{op-tical net-works semi-conduc-tor}


\begin{document}

\title{Project Smart Trigger \\ Design Documentation \\ for ECE/CS 3992}

\author{Andrew Koebbe \textit{Student ID: u1450645}, Stephen Hofmann \textit{Student ID: u1028605},\\ Thomas Stay \textit{Student ID: u1249316}, Derek Jenkins \textit{Student ID: u1506428}}%

\maketitle

\begin{abstract}
This document describes the design of the Smart Trigger, a modular RFID-based
access-control system for the AR-15 platform. A passive 125kHz RFID tag worn
by the authorized user is read by an antenna embedded in a custom 3D-printed
grip. A Raspberry Pi Pico microcontroller validates the tag against an
authorized list and drives a micro stepper motor that moves a blocking member
out of the fire control group. The document covers the system architecture, the
hardware components and their integration onto a custom PCB, the CAD modeling
of the grip and actuator mounts, and the firmware processes that tie the system
together. It closes with a look at how the design could grow through an
off-device logging server and web application, and what that complexity would
cost in power, security, and development effort.
\end{abstract}

\IEEEpeerreviewmaketitle

\section{Introduction}

\section{System Architecture}
\subsection{Overview}
The Smart Trigger is a self-contained, battery-powered embedded system that
adds an electronic authorization step to the mechanical operation of an AR-15.
Functionally, it is a closed loop with three stages: Sense, Decide, Act.

\subsection{Subsystems}

\subsection{Design Requirements}

\subsection{Key Architectural Decisions}


\section{Hardware Components and Integration}
\subsection{Mechanical Interaction}
The mechanical design has three parts:
\begin{enumerate}
    \item a \textbf{custom grip} that replaces the factory AR-15 grip and
    houses the electronics, battery, and RFID antenna
    \item an \textbf{actuator mount} that locates the micro stepper in the fire
    control pocket of the lower receiver
    \item a \textbf{blocking member} driven by the motor that physically
    prevents the fire control group from releasing the hammer
\end{enumerate}

\subsection{Component Selection}


\subsection{Main Board (PCB)}

\subsection{RFID Board (RDM6300)}

\subsection{Micro Stepper Motor & Driver Board (SAA1042)}

\subsection{Integration and Test Plan}


\section{CAD Modeling and Design}
\label{sec:cad}

\subsection{Scope of the Mechanical Design}

\subsection{Custom Grip}

\subsection{Actuator Mount and Blocking Member}

\subsection{Materials and Print Parameters}

\subsection{Design Iteration}



\section{Software Processes}
\label{sec:software}

\subsection{Development Environment}

\subsection{Firmware Architecture}

\subsection{Control State Machine}

\subsection{RFID Processing}

\subsection{Actuator Control}

\subsection{Software Testing}

\subsection{Expansion: Logging Server and Web Application}


\section{Conclusion}

\clearpage

\bibliographystyle{IEEEtran}
\bibliography{bibtex/bib/design}

\end{document}
